\section*{第 \ref{ch:b1:second-law-entropy} 章——热力学第二定律：熵}

\begin{solution}{exo:b1:second-law-entropy:1}
$\Delta S = mL/T$：熔化 $334000/273 = \qty{1.22}{kJ/K}$；沸腾 $2260000/
373 = \qty{6.06}{kJ/K}$。汽化使分子彼此
完全脱离——微观无序度大幅增加，
每\omterm{def:b1:kinetic-theory:temperature}{开尔文}对应的热量也多得多。
\end{solution}

\begin{solution}{exo:b1:second-law-entropy:2}
$\Delta S = mc\ln(T_2/T_1) = 4180\ln(353/293) = \qty{0.78}{kJ/K}$——状态
函数：与途径无关（\emph{产生}的熵则不然）。
\end{solution}

\begin{solution}{exo:b1:second-law-entropy:3}
$\Delta S = nR\ln2 = \qty{5.8}{J/K}$。可逆\omterm{prop:b1:first-law:transformations}{等温过程}：$Q = nRT\ln2$
由温度为 $T$ 的\omterm{def:b1:first-law:heat}{热源}提供，$S_{\text{交换}} = nR\ln2$，产生的熵为 $0$。焦耳实验：
$Q = 0$，$S_{\text{交换}} = 0$，产生的熵为 \qty{5.8}{J/K}。
\end{solution}

\begin{solution}{exo:b1:second-law-entropy:4}
$S = nC_{V,m}\ln T + nR\ln V$ 为常量：$T^{C_{V,m}}V^R$ 为常量，即\
$TV^{R/C_{V,m}} = TV^{\gamma - 1}$ 为常量。突然压缩时：$Q = 0$，故
$S_{\text{交换}} = 0$，且 $\Delta S = S_{\text{产生}}$；$T$：$300 \to
643$ K，$V$：$24.9 \to 10.7$ L：$\Delta S = 20.8\ln(643/300) + 8.314\ln(10.7/
24.9) = 15.9 - 7.0 = \qty{+8.9}{J/K}$，为产生的熵。
\end{solution}

\begin{solution}{exo:b1:second-law-entropy:5}
物块：$450\ln(293/373) = \qty{-109}{J/K}$；湖水：$+36000/293 = \qty{+123}{J/K}$；
宇宙：$+\qty{14}{J/K}$。
\end{solution}

\begin{solution}{exo:b1:second-law-entropy:6}
$T_f = \qty{350}{K}$；$S = 4180[\ln(350/400) + \ln(350/300)] = 4180 \times
0.0206 = \qty{86}{J/K}$，全部为产生的熵。$(T_1 + T_2)^2 - 4T_1T_2 = (T_1 - T_2)^2
\geq 0$，故对数的真数 $\geq 1$。
\end{solution}

\begin{solution}{exo:b1:second-law-entropy:7}
$S = Q(1/T_c - 1/T_h)$：$1000(1/280 - 1/300) = \qty{0.24}{J/K}$；从炉子到
房间：$1000(1/300 - 1/1000) = \qty{2.3}{J/K}$，约为前者的十倍：热流跨越的温差越大，
过程就越不可逆——炉子到房间的温降，
正是高品位能量被白白浪费之处。
\end{solution}

\begin{solution}{exo:b1:second-law-entropy:8}
循环：$\Delta S = 0 = Q/T_0 + S_{\text{产生}}$，且 $S_{\text{产生}} \geq 0$，
故 $Q \leq 0$，$W = -Q \geq 0$：机器无法对外做功。
海洋是单一的\omterm{def:b1:first-law:heat}{热源}：若没有更冷的物体来接纳熵，
其热量就无法转化为功。
\end{solution}

\begin{solution}{exo:b1:second-law-entropy:9}
每种气体由 $V$ 膨胀到 $2V$（温度 $T$ 保持不变）：$\Delta S = 2 \times R\ln2
= \qty{11.5}{J/K}$，产生的熵（没有热量交换，箱子是孤立的）。
两半氮气：抽去隔板并不改变宏观状态——
$\Delta S = 0$（吉布斯佯谬：相同分子不算“混合”）。
\end{solution}

\begin{solution}{exo:b1:second-law-entropy:10}
$Q = \qty{3.6e5}{J}$ 传入 \qty{293}{K} 的房间：$S = \qty{1230}{J/K}$
产生的熵——一部分在灯丝中（电功转化为热），
另一部分在该热量从灼热的灯丝流向较冷房间的过程中；电
功本身不携带熵。
\end{solution}

\begin{solution}{exo:b1:second-law-entropy:11}
$\Delta S_{\text{物体}} = C\ln(T_0/T_1)$，各种情形下均成立。(a) 交换的热量
为 $C(T_0 - T_1)/T_0$，产生的熵为 $C[T_1/T_0 - 1 - \ln(T_1/T_0)]$。(b) 分两步，
每步的比值 $r = \sqrt{T_1/T_0}$，产生的熵为 $2C[r - 1 - \ln r]$，更小
（取 $T_1/T_0 = 2$ 时为 $0.193C$，对比 $0.307C$）。(c) 无穷多步：
每步产生的熵为 $C[(1 + \epsilon) - 1 - \ln(1 + \epsilon)] \approx C\epsilon^2/2$，
当 $\epsilon \to 0$ 时求和趋于零：可逆冷却需要在每个中间温度处都有
\omterm{def:b1:first-law:heat}{热源}——热量绝不能跨越有限的温差。
\end{solution}

\begin{solution}{exo:b1:second-law-entropy:12}
$2^{-10} = 10^{-3}$；$2^{-100} = \num{8e-31}$；$2^{-10^{22}} = 10^{-3\times10^{21}}$。
$k_B\ln2^N = Nk_B\ln2 = nR\ln2$。“不可能”严格地说意味着概率为
零；这里的概率是 $10^{-3\times10^{21}}$——并非零，而是一个小到
任何差异都永远无法被观测的数：不可逆性是在
$10^{23}$ 量级上的统计规律。
\end{solution}

\begin{solution}{pb:b1:second-law-entropy:1}
\textbf{1.} $Q = mc(T_0 - T_1) = 0.25 \times 4180 \times (-70) = \qty{-73}{kJ}$
对咖啡而言：房间获得 \qty{73}{kJ}，$\Delta S_{\text{房间}} = 73150/293
= \qty{+250}{J/K}$。

\textbf{2.} $mc\ln(T_0/T_1) = 1045\ln(293/363) = \qty{-224}{J/K}$.

\textbf{3.} $S_{\text{产生}} = -224 - (-250) = \qty{+26}{J/K} > 0$.

\textbf{4.} $\Delta S = 1045\ln(330/363) = \qty{-99.6}{J/K}$; $Q = 1045 \times
(-33) = \qty{-34.5}{kJ}$，$S_{\text{交换}} = \qty{-117.7}{J/K}$：目前产生的熵
为 \qty{18}{J/K}——大部分不可逆性出现在
咖啡最烫的时候。

\textbf{5.} $\Delta S - Q/T_0 = mc[\ln(T_0/T_1) - (T_0 - T_1)/T_0] = mc[x - 1 -
\ln x]$，其中 $x = T_1/T_0$；$f'(x) = 1 - 1/x$ 在 $x = 1$ 处为零，
此时 $f = 0$，且 $f'' = 1/x^2 > 0$：极小值为 $0$。

\textbf{6.} 否：初态和终态相同，热量相同，
房间是同一个\omterm{def:b1:first-law:heat}{热源}——产生的熵由
状态决定，而与快慢无关。

\textbf{7.} $T_f = \qty{350}{K}$；产生的熵为 $4180\ln(350^2/120000) = \qty{86}{J/K}$。

\textbf{8.} 可逆：总熵守恒：$mc\ln(T_f'/T_1) + mc\ln(T_f'/
T_2) = 0$，故 $T_f'^2 = T_1T_2$。

\textbf{9.} $T_f' = \sqrt{120000} = \qty{346.4}{K}$；水失去的热量为 $mc[(T_1 -
T_f') + (T_2 - T_f')] = 4180 \times (53.6 - 46.4) = \qty{30}{kJ}$，
以功的形式输出。

\textbf{10.} 热水：$4180\ln(346.4/400) = \qty{-601}{J/K}$；冷水：$4180\ln(346.4/
300) = \qty{+601}{J/K}$：总和为零，符合\omterm{def:b1:first-law:transformation}{可逆过程}的要求。

\textbf{11.} $T_0S_{\text{产生}} = 300 \times 86 = \qty{26}{kJ}$，接近
可回收的 \qty{30}{kJ}：产生的熵乘以参考温度，
即为\omterm{rem:b1:first-law:irreversible}{不可逆过程}所浪费的功。

\textbf{12.} 需要一台工作于两杯温水之间的热机，而水温在不断变化——
对于大型机器也只有几十千焦：不值得去做，
但热力学第二定律表明，这份功
确实存在。

\textbf{13.} 从内部吸收 $Q_c/T_c = 40/278 = \qty{0.144}{W/K}$；
传给房间 $Q_0/T_0 = (40 + W)/293$；一个循环内的平衡：
$0 = 0.144 - (40 + W)/293 + S_c$。

\textbf{14.} $S_c \geq 0$：$(40 + W)/293 \geq 0.144$，$W \geq 42.2 - 40 =
\qty{2.2}{W}$；性能系数 $\leq 40/2.2 = 18$（$= T_c/(T_0 - T_c)$）。

\textbf{15.} $\mathrm{COP}_{\max} = T_c/(T_0 - T_c)$：对于冷冻柜，$255/38 = 6.7$
——箱内越冷，每移出一焦耳热量所付出的代价就越大。

\textbf{16.} 取 $W = 15$：$S_c = (55/293) - 0.144 = \qty{0.044}{W/K}$，
即每秒产生的熵。

\textbf{17.} $W \geq 200 \times 15/278 = \qty{10.8}{W}$；机器将
\qty{200}{W} 加上自身的功一并排入房间，而内部保持寒冷：
房间的净得热为 $W$——敞开的冰箱就是一台加热器。

\textbf{18.} $W = 0$ 将给出 $S_c = Q_c(1/T_0 - 1/T_c) < 0$：被禁止——
这正是克劳修斯的表述。

\textbf{19.} 初态：$1$（全部在左侧）；终态：$2^N$；$\Delta S = k_B\ln2^N =
Nk_B\ln2$。

\textbf{20.} $Nk_B = nR$：$\Delta S = nR\ln2$，即\omterm{rem:b1:first-law:irreversible}{焦耳膨胀}的结果。

\textbf{21.} 每个分子对应一个比特：共 $N$ 个比特；且 $\Delta S/(k_B\ln2) = N$——
熵度量的是确定这些分子
位置所需的信息量。

\textbf{22.} $2^{-50} = \num{9e-16}$；对于一\omterm{def:b1:kinetic-theory:scales}{摩尔}，$2^{-6\times10^{23}} =
10^{-1.8\times10^{23}}$。

\textbf{23.} 无热量交换（$Q = 0$）：交换的熵为 $0$，产生的熵为 $nR\ln2$——
最纯粹的情形：无功，无热，只有气体的扩散。

\textbf{24.} 接受的功为 $W = nRT\ln2$，给予房间的热量为 $Q = -nRT\ln2$，
$S_{\text{交换}} = -nR\ln2$（气体的熵回到初始值，过程可逆）：
撤销自由膨胀需要做功 $nRT\ln2$，
并向房间排出 $nR\ln2$ 的熵——宇宙永久地保留了
一次性产生的 $nR\ln2$。

\textbf{25.} 现实中没有任何过程能守恒熵——
熵守恒只存在于理想的可逆极限中；定律
禁止每一个会破坏熵的过程（热量逆向流动、从单一\omterm{def:b1:first-law:heat}{热源}吸热做功、无代价地分离混合物）；
每产生一个比特，就损失一份功，其数值为 $T_0S_{\text{产生}}$。
\end{solution}
